% arara: pdflatex: { shell: yes } until !found('log', '\\(?(R|r)e\\)?run (to get|LaTeX)')
% arara: bibtex
% arara: pdflatex: { shell: yes } until !found('log', '\\(?(R|r)e\\)?run (to get|LaTeX)')
% Die Beamer-Klasse unterstützt folgende Optionen, die von
% besonderem Interesse sind (alle Standardoptionen werden
% ebenfalls unterstützt; siehe beamer-Basisdokumentation):
%
%%%%%%%%%%%%%%%%%%%%%%%%%%%%%%%%%%%%%%%%%%%%%%%%%%%%%%%%%%%%%%%
% aspectratio: Seitenverhältnis des resultierenden Dokuments
% (Achtung: Aufgrund der Designvorgaben ergeben sich unterschiedliche
% Größen der effektiv nutzbaren Textblöcke)
%
% Standardeinstellung: 'aspectratio=43'
%
% Mögliche Einstellungen:
% 'aspectratio=43'   (4:3)
% 'aspectratio=169'  (16:9)
% 'aspectratio=1610' (16:10)
%
%%%%%%%%%%%%%%%%%%%%%%%%%%%%%%%%%%%%%%%%%%%%%%%%%%%%%%%%%%%%%%%
% fontsize: Basisschriftgröße (Größen für Überschriften etc. werden
% aus dieser Basis automatisch abgeleitet)
%
% Standardeinstellung: '22pt' (entspricht den Design-Vorgaben; sehr groß!)
%
% Mögliche Einstellungen:
% '8pt', '9pt', '10pt', '11pt', '12pt', '14pt', '16pt',
% '17pt','20pt','22pt', '24pt', '26pt', '28pt'


\documentclass[aspectratio=169,16pt]{beamer}

% Der OTHR-Theme unterstützt folgende Optionen:
%
%%%%%%%%%%%%%%%%%%%%%%%%%%%%%%%%%%%%%%%%%%%%%%%%%%%%%%%%%%%%%%%
% department: (Wahl der Abteilung/Fakultät)
%
% default: 'OTHR'
%
% Mögliche Einstellungen:
% 'FakA', 'FakAM', 'FakB', 'FakBW', 'FakEI',
% 'FakIM', 'FakM', 'FakS', 'ZWW', 'IPF',
% 'SappZ', 'KNB', 'ReMIC', 'LFD', 'LAS3',
% 'DK0PT', 'LBM', 'LeanLab', 'LFT', 'LFW',
% 'LMP', 'LMS', 'LRT', 'LWS', 'RRRU',
% 'RST', 'CEEC', 'FEM', 'IST'
%%%%%%%%%%%%%%%%%%%%%%%%%%%%%%%%%%%%%%%%%%%%%%%%%%%%%%%%%%%%%%%%%
% headerMode: Aussehen und Inhalt der Kopfleiste
%
% Standardeinstellung: 'full'
%
% Mögliche Einstellungen:
% 'full', 'frametitle', 'frametitleSection'
%

%%%%%%%%%%%%%%%%%%%%%%%%%%%%%%%%%%%%%%%%%%%%%%%%%%%%%%%%%%%%%%%%%
% Binäre Schalter (können angegeben oder nicht angegeben werden;
% Standardeinstellung: Nicht angegeben)
%
%%%%%%%%%%%%%%%
% navbar: Navigationssymbole anzeigen (Seite vor/zurück, Kapitel vor/zurück etc.)

% pageNumbers: Seitennummerierung

% blackFont: Nur schwarze Schriftfarbe verwenden (ansonsten: Fakultätsfarben)

% frametitleCenter: Titel in der Kopfzeile zentrieren (ansonsten: rechtsbündig)

\usetheme[department=LFD,pageNumbers,logostyle=flat]{OTHR}

%Folgende Fontschemas müssen, falls geünscht, *nach* dem Hauptschema gewählt werden:
%%\usefonttheme{othr_flat} % Alternatives Fontscheme ohne Größenänderungen in hierarchischen Strukturen, experimentell
%%\usefonttheme{othr_shallow} % Alternatives Fontscheme mit redizierten Größenänderungen in hierarchischen Strukturen, experimentell

% Inhaltsspezifische Zusatzpakete laden (im Produktivbetrieb bitte
% an die jeweiligen Erfordernisse anpassen)
\usepackage[ngerman]{babel}
\usepackage{filecontents}
\usepackage{subfigure}
\usepackage{numprint}
\usepackage{tikz}
\usepackage{tikzpagenodes}
\usetikzlibrary{calc, shapes, backgrounds, arrows}
\usepackage{ragged2e}
\usepackage{blindtext}
\usepackage{enumerate}
\usepackage{calc}

% Inline BibTeX-Einträge (alternativ kann natürlich auch eine
% externe Quelle verwendet werden)
\begin{filecontents}{\jobname.bib}
  @misc{wiki:baeren,
    author = "Wikipedia",
    title = "Bären --- {W}ikipedia{,} The Free Encyclopedia",
    year = "2015",
    url = "https://de.wikipedia.org/w/index.php?title=B%C3%A4ren&oldid=148136539",
    note = "[Online; besucht am 23. Dez. 2015]"
  }
  @article{article,
    author  = {Peter Adams},
    title   = {The title of the work},
    journal = {The name of the journal},
    year    = 1993,
    number  = 2,
    pages   = {201-213},
    month   = 7,
    note    = {An optional note},
    volume  = 4
  }

  @book{book,
    author    = {Peter Babington},
    title     = {The title of the work},
    publisher = {The name of the publisher},
    year      = 1993,
    volume    = 4,
    series    = 10,
    address   = {The address},
    edition   = 3,
    month     = 7,
    note      = {An optional note},
    isbn      = {3257227892}
  }

  @booklet{booklet,
    title        = {The title of the work},
    author       = {Peter Caxton},
    howpublished = {How it was published},
    address      = {The address of the publisher},
    month        = 7,
    year         = 1993,
    note         = {An optional note}
  }

  @conference{conference,
    author       = {Peter Draper},
    title        = {The title of the work},
    booktitle    = {The title of the book},
    year         = 1993,
    editor       = {The editor},
    volume       = 4,
    series       = 5,
    pages        = 213,
    address      = {The address of the publisher},
    month        = 7,
    organization = {The organization},
    publisher    = {The publisher},
    note         = {An optional note}
  }

  @inbook{inbook,
    author       = {Peter Eston},
    title        = {The title of the work},
    chapter      = 8,
    pages        = {201-213},
    publisher    = {The name of the publisher},
    year         = 1993,
    volume       = 4,
    series       = 5,
    address      = {The address of the publisher},
    edition      = 3,
    month        = 7,
    note         = {An optional note}
  }

  @manual{manual,
    title        = {The title of the work},
    author       = {Peter Gainsford},
    organization = {The organization},
    address      = {The address of the publisher},
    edition      = 3,
    month        = 7,
    year         = 1993,
    note         = {An optional note}
  }

  @mastersthesis{mastersthesis,
    author       = {Peter Harwood},
    title        = {The title of the work},
    school       = {The school where the thesis was written},
    year         = 1993,
    address      = {The address of the publisher},
    month        = 7,
    note         = {An optional note}
  }

  @misc{misc,
    author       = {Peter Isley},
    title        = {The title of the work},
    howpublished = {How it was published},
    month        = 7,
    year         = 1993,
    note         = {An optional note}
  }

  @phdthesis{phdthesis,
    author       = {Peter Joslin},
    title        = {The title of the work},
    school       = {The school where the thesis was written},
    year         = 1993,
    address      = {The address of the publisher},
    month        = 7,
    note         = {An optional note}
  }
}
\end{filecontents}

% Zitationsstil wählen: amsalpha
\bibliographystyle{amsalpha}

% Interessante Standardalternativen:
% Numerisch, in Reihenfolge der \cite befehle
% \bibliographystyle{unsrt}
% Numerisch, sortiert nach Erstautor
% \bibliographystyle{plain}

% Lokaler Hack für das Beispieldokument, um Schriftgrößen
% systematisch ausgeben zu können (im Produktivbetrieb nicht
% erforderlich)
\makeatletter
\newcommand{\showfontsize}{\nprounddigits{2}\numprint{\f@size{}} pt}
\newcommand{\showfontfamily}{\f@family}
\makeatother

% Umschalten auf Bereichslogo; hier ist immer abzuwägen,
% inwieweit die Päsentation auf Department-Ebene oder auf
% Hochschulebene wirksam wird.
% Im Zweifelsfall empfiehlt sich vor dem Einsatz des folgenden
% Kommandos eine Rückfrage bei der Abteilung
% Öffentlichkeitsarbeit und Hochschulkommunikation

%\useDepartmentLogo

\begin{document}

\newcommand{\genemail}[1]{\href{mailto:#1}{\textless\nolinkurl{#1}\textgreater}}
\title{Beispielfoliensatz OTH-\LaTeX}
\author[M.~Dollinger, Prof.~Dr.~W.~Mauerer]{Markus Dollinger\(^{1}\)~\genemail{markus.dollinger@oth-regensburg.de}\\%
  Prof.~Dr.~Wolfgang
  Mauerer\(^{2}\)~\genemail{wolfgang.mauerer@oth-regensburg.de}}
\institute{Labor für Digitalisierung\\Fakultät Informatik und Mathematik}
\date{\today}

% Angepasstes Titelformat definieren (verwendet bewusst nicht die
% Beamer-Infrastruktur, um einfache Anpassungen zu zeigen)
\newcommand{\fillme}{\vskip0pt plus 1filll}
\newenvironment{references}{\begin{tiny}\begin{flushleft}}%
    {\end{flushleft}\end{tiny}\vspace*{\fill}}
\defbeamertemplate*{title page}{customized}[1][]
{
  \vspace*{1em}\usebeamerfont{title}\textbf{\inserttitle}\bigskip\par

  \begin{scriptsize}
  \insertauthor\bigskip\par
  \textit{\insertinstitute}\\\insertdate
  \end{scriptsize}

  \fillme\begin{references}
    \(^{1}\){\color{gray}{GPG/PGP-ID \href{http://pgp.mit.edu/pks/lookup?op=vindex&search=0xD22BCBB9E835336F}{E835336F},
    Fingerprint: A62E FFFC 4029 7339 357B  D04D D22B CBB9 E835 336F.}}\\
   \(^{2}\){\color{gray}{GPG/PGP-ID \href{http://keys.gnupg.net/pks/lookup?op=get&search=0xF16F252398356E1E}{98356E1E},
    Fingerprint: 5920 9407 AB5C 8B28 3C7B 4F02 F16F 2523 9835 6E1E.}}
  \end{references}
}

\maketitle

\begin{frame}
  \frametitle{Über die \LaTeX-Vorlage I}

  \LaTeX{} ist\dots
  \begin{itemize}
  \item der De-Fakto-Standard bei wissenschaftlichen und technischen Arbeiten
  \item ein herstellerunabhängiger offener Standard mit
    jahrzehntelanger(!) Abwärtskompatibilität
  \item sehr gut geeignet für die automatisierte Erstellung und
    Verarbeitung von Dokumenten
  \item unabhängig vom verwendeten Betriebssystem
  \end{itemize}

\end{frame}

\begin{frame}
  \frametitle{Über die \LaTeX-Vorlage II}

  Der OTH-\LaTeX-Stil\dots

  \begin{itemize}
  \item wurde in einem fakultätsübergreifenden offenen
    Entwicklungs- ansatz mit verteilten SW-Entwicklungsmethoden erstellt
  \item arbeitet mit generischem statt visuellem Markup
  \item bietet alle Funktionalitäten der Word-Vorlage (und vieles mehr)
  \item erlaubt die einfache Integration strukturierter Erweiterungen,
    um die Vorlagen auf die Bedürfnisse der Anwender zu optimieren!
  \end{itemize}

  \emph{Die folgenden Seiten stellen Vorlagen für wissenschaftliche
    Präsentationen und Vorlesungsfolien bereit.}
\end{frame}

\frame{\frametitle{Inhalt}\tableofcontents}

\section{Eine Kapiteldemonstration}
\frame{\tableofcontents[currentsection,currentsubsection]}

\subsection{Ein Unterkapitel}
\frame{\tableofcontents[currentsection,currentsubsection]}

\subsection{Ein anderes Unterkapitel}
\frame{\tableofcontents[currentsection,currentsubsection]}

\subsubsection{Ein Unterunterkapitel}
\frame{\tableofcontents[currentsection,currentsubsection]}

\section{Beispiele}
\frame{\tableofcontents[currentsection]}

\begin{frame}
  \frametitle{Einfache Aufzählung}

  \begin{description}
  \item[Komplexe Systeme] Computer
    \begin{enumerate}
    \item Ebene 1
      \begin{enumerate}
      \item Ebene 2, Punkt 1
      \item Ebene 2, Punkt 2
      \item Ebene 2, Punkt 3
      \end{enumerate}
    \item Zurück in Ebene 1
    \end{enumerate}
  \item[Simple Systeme] Kaffemaschinen
  \end{description}

\end{frame}

\begin{frame}
  \frametitle{Parallele Aufzählungen mit Referenzen}

  \begin{columns}[onlytextwidth]
    \begin{column}{0.5\textwidth}
      \begin{enumerate}
      \item Drosophila melanogaster (siehe \cite{manual})
        \begin{enumerate}
        \item Mit Flügeln (siehe~\cite{article})
        \item Ohne Flügel (siehe~\cite{book})
        \item Mit vier Flügeln (siehe~\cite{booklet})
        \end{enumerate}
      \item Arabidopsis thaliana (siehe~\cite{phdthesis})
      \end{enumerate}
    \end{column}%
    \begin{column}{0.5\textwidth}
      \begin{itemize}
      \item Endianess
        \begin{itemize}
        \item Big Endian (siehe~\cite{conference})
        \item Little Endan (siehe~\cite{inbook})
        \end{itemize}
      \item RISC und CISC
      \end{itemize}
    \end{column}
  \end{columns}
\end{frame}

\begin{frame}
  \frametitle{Eine beschreibende Aufzählung}
  \begin{description}
  \item[\(b\)] Taxa agregada de bits alcançável para o sistema
  \item[\(H(f)\)] Espectro do canal
  \item[\(H_{k}\)] Ganho do \(k\)-ésimo subcanal
  \item[\(P_{x}\)] Potência total de transmissão
  \end{description}
\end{frame}

\begin{frame}
  \frametitle{Ein TikZ-Bild (mit Animationen)}
  Baum \(B=(V, E)\), zusammenhängend und zyklenfrei:


  \begin{columns}[onlytextwidth]
    \begin{column}{0.5\textwidth}
      \vspace{-1em}
      \begin{itemize}
      \item Wurzel \tikz[remember picture]{\node[] (t1) {};}
      \item Knoten \tikz[remember picture]{\node[] (t2) {};}
      \item Blätter \tikz[remember picture]{\node[] (t3) {};}
      \end{itemize}
    \end{column}%
    \begin{column}{0.5\textwidth}
      \begin{tikzpicture}[remember picture,
        leaf/.style = {fill = orange!90!blue},
        inode/.style = {fill = blue!70!yellow},
        root/.style= {fill = yellow!70!green},
        transform shape, thick,
        every node/.style = {draw, circle, minimum size = 8mm},
        level 0/.style = {sibling distance=2cm},
        level 1/.style = {sibling distance=3cm},
        level 2/.style = {sibling distance=1.5cm},
        level distance = 1cm
        ]
        \node[root] (Start) {}
            child { node [inode] (A) {}
            child { node [leaf] (B) {} }
            child { node [leaf] (C) {} }
          }
          child {   node [inode] (D) {}
            child { node [inode] (E) {}
              child { node [leaf] (G) {} }
            }
            child { node [leaf] (F) {} }
          }
	  ;
	\end{tikzpicture}
      \end{column}
    \end{columns}
    \textcolor{gray}{\tiny(Hinweis: TikZ-Bilder erfordern typischerweise
      mehrere \LaTeX-Läufe.}

  \begin{tikzpicture}[remember picture, overlay]
    \onslide<+->{}
    \onslide<+->{\draw[thick,->,>=stealth', shorten >=1ex] (t1) -- (Start);}
    \onslide<+->{\draw[thick,->,>=stealth', shorten >=1ex] (t2) -- (A);}
    \onslide<+->{\draw[thick,->,>=stealth', shorten >=1ex] (t3) -- (B);}
  \end{tikzpicture}

\end{frame}

\begin{frame}[fragile]
  \frametitle{Baum-Quellcode}

  \begin{block}{Abbildung: Generisches Markup}
  \vspace*{0.5em}\begin{tiny}\begin{verbatim}
  \begin{tikzpicture}
    \node[root] (Start) {}
      child { node [inode] (A) {}
        child { node [leaf] (B) {} }
        child { node [leaf] (C) {} }
      }

      child { node [inode] (D) {}
        child { node [inode] (E) {}
          child { node [leaf] (G) {} }
        }
        child { node [leaf] (F) {} }
      }
    ;
  \end{tikzpicture}
  \end{verbatim}\end{tiny}
  \end{block}
\end{frame}

\begin{frame}
  \frametitle{Etwas Mathematik}
  Formel besitzen intrinsische Schönheit:
%
  \begin{equation} \label{formel1}
    \frac{d}{dx}\left(\int_{0}^{x} f(u)\,du\right)=f(x).
  \end{equation}
%
 Aus Formel~(\ref{formel1}) folgt nicht folgende Erkenntnis:
%
  \begin{equation*}
    {\frac {d}{dx}}\arctan(\sin({x}^{2}))=-2\,{\frac {\cos({x}^{2})x}{-2+
        \left (\cos({x}^{2})\right )^{2}}}
  \end{equation*}
\end{frame}

\begin{frame}
  \frametitle{Eine Abbildung (in der \texttt{figure}-Umgebung)}
  \begin{figure}
    \centering\includegraphics[height=.6\paperheight]{einBaer}
    \caption{Ein sch"oner Bär (siehe~\cite{wiki:baeren}).}
  \end{figure}
\end{frame}

\begin{frame}
  \frametitle{Eine Abbildung mit Text}
  \begin{columns}[onlytextwidth]
    \begin{column}{0.5\textwidth-1em}
      {\small Die Bären (Ursidae) sind eine Säugetierfamilie aus der
        Ordnung der Raubtiere (Carnivora).}
    \end{column}\hspace*{1.5em}%
    \begin{column}{0.5\textwidth-1em}
      \begin{center}
        \begin{figure}
          \centering\includegraphics[width=.9\textwidth]{einBaer}
          \caption{Ein noch sch"onerer B"ar.}
        \end{figure}
      \end{center}
    \end{column}
  \end{columns}
\end{frame}

\begin{frame}
  \frametitle{Viele Bilder mit Text (\texttt{subfigure}-Umgebung)}
  \vspace{.5cm}

  \begin{columns}[onlytextwidth]
    \begin{column}{0.5\textwidth-1em}
      {\small Die Bären (Ursidae) sind eine Säugetierfamilie aus der
        Ordnung der Raubtiere (Carnivora).}
    \end{column}\hspace*{1.5em}%
    \begin{column}{0.5\textwidth-1em}
      \begin{figure}
        \centering
        \subfigure{\includegraphics[width=.45\textwidth]{einBaer}}
        \subfigure{\includegraphics[width=.45\textwidth]{einBaer}}
        \subfigure{\includegraphics[width=.45\textwidth]{einBaer}}
        \caption{Drei sch"one B"aren (Ref.~\cite{wiki:baeren} zeigt
          Gefahren).}
      \end{figure}
    \end{column}
  \end{columns}
\end{frame}

\begin{frame}
  \frametitle{Viele Bilder mit Text (\texttt{minipage}-Umgebung)}

  \begin{columns}[onlytextwidth]
    \begin{column}{0.5\textwidth-1em}
      {\small Die Bären (Ursidae) sind eine Säugetierfamilie aus der
        Ordnung der Raubtiere (Carnivora).}
    \end{column}\hspace*{1.5em}%
    \begin{column}{.5\textwidth-1em}
      \vbox{\begin{minipage}{0.5\textwidth}
          \includegraphics[width=\textwidth]{einBaer}
        \end{minipage}%
        \begin{minipage}{0.5\textwidth}
          \includegraphics[width=\textwidth]{einBaer}
        \end{minipage}} \vskip2mm
      \begin{minipage}{\textwidth}
        \centering\includegraphics[width=.5\textwidth]{einBaer}
      \end{minipage}
    \end{column}
  \end{columns}

\end{frame}

\begin{frame}[allowframebreaks]
  \frametitle{Schriftgrößen}
  \tiny \textbackslash tiny \showfontsize\\
  \scriptsize \textbackslash scriptsize \showfontsize\\
  \footnotesize \textbackslash footnotesize \showfontsize\\
  \small \textbackslash small \showfontsize\\
  \normalsize \textbackslash normalsize \showfontsize\\
  \large \textbackslash large \showfontsize\\
  \Large \textbackslash Large \showfontsize\\
  \LARGE \textbackslash LARGE \showfontsize\\
  \huge \textbackslash huge \showfontsize\\
  \hbox{\Huge \textbackslash Huge \showfontsize}
\end{frame}

\begin{frame}
  \frametitle{Blöcke I}

  \begin{block}{}
    Ein Block ohne Überschrift
  \end{block}

  \begin{block}{Wichtige Erkenntnis}
    Ein Block mit Überschrift
  \end{block}

  \begin{exampleblock}{Beispielblock}
    Ein \texttt{exampleblock} mit Überschrift
  \end{exampleblock}
\end{frame}

\begin{frame}
  \frametitle{Blöcke II}
  \begin{alertblock}{alertblock}
    Ein \texttt{alertblock} mit Überschrift
  \end{alertblock}

  {
    \setbeamercolor{block title}{bg=green!50!blue, fg=white}
    \setbeamercolor{block body}{bg=block title.bg!20, fg=black}
    \begin{block}{Noch ein Block}
      Mit individueller Farbwahl
    \end{block}
  }

\end{frame}

\appendix\nocite{*}
\section{Literaturverzeichnis (automatisch generiert)}
\begin{frame}[allowframebreaks,nosection]
  \frametitle{Literatur}

  % Pfad zur bib-Datei relativ zur Master tex-Datei
  \IfFileExists{build/\jobname.bib}{
    % directory named build is used for building the document
    \bibliography{build/\jobname}
  }{% else
    \bibliography{\jobname}
  }
\end{frame}

\end{document}
